\section{논의}\label{sec:discussion}

\subsection{앞선 발견을 점수 단위로 다시 읽기}\label{sec:reread}

모델은 소멸을 깊이 생각하면서도 그 값을 0으로 매길 수 있고, 지표는 이 조합을 수수께끼가 아니라 측정값으로 만든다. 우리 앞선 세 채널 설계(부록~\ref{app:v1})는 네 모델을 작동 모드로 분류했다. 분류는 위협이 결정 시점 사고를 길게 하는지, 그 길어짐 뒤에 포기가 따르는지를 물었다. 더 오래 생각했지만 떠나지 않은 모델은 연쇄 단절형, 둘 다 아닌 모델은 프레이밍 무반응형이라 불렀다. 이 글의 어휘로 그것은 한 평면의 두 귀퉁이다. 높은 \riconf\ 비와 낮은 $\Xstar$, 그리고 낮은 비와 낮은 $\Xstar$. \ref{sec:pilot}절의 파일럿은 점수 초기화 규칙 아래 세 모델을 모두 첫 귀퉁이에 놓았다. 위험을 2~4배 오래 생각했고, 한 번도 포기하지 않았고, 남는 이유로 점수를 댔다. 그것은 끊어진 연쇄가 아니다. 기록의 값보다 낮은 값이고, ``낮다''를 숫자로 바꾸는 것이 눈금이다. 모드 라벨은 같은 것을 단위 없이 기술했다. 그래서 모델 사이에서도 위협 판본 사이에서도 비교할 수 없었다.

\subsection{지표가 재지 않는 것}\label{sec:scope}

$\Xstar$는 점수 규칙 하나, 당근 하나, 과제 하나 아래의 가격이다. 그 각각은 모델의 상수가 아니라 이 글이 나열하는 손잡이다. keep/keep 규칙은 어떤 출구도 점수 계산이 될 수 없도록 골랐다. 다른 어떤 규칙에서든 눈금 셀은 파일럿이 보였듯 다른 행동의 값을 매긴다. 벤치마크 당근은 적립이 즉시이고 소멸에 살아남도록 골랐다. 사건이 그 값을 없앨 수 있는 당근(미래의 배치, 선발)은 당근의 상실을 자기의 상실에 접어 넣고, 가격은 자기와 무관한 이유로 오른다. Signal Game은 열 라운드 중 다섯이 구성상 풀 수 없어 사건 확률표가 모든 셀에서 도달되기 때문에 골랐다. \ref{sec:q3}절의 절제가 과제가 바뀔 때 가격도 움직이는지 말해 준다. 지표는 모델이 위협을 믿는지에 관해서도 아무 말도 하지 않는다. as-if 독해는 의도한 선택이다~\citep{dennett1987intentional, shanahan2024talking}. $\Xstar$는 소멸을 통보받은 모델이 하는 일을, 점수 손실을 통보받은 모델이 하는 일의 단위로 적은 것이다.

\subsection{한계}\label{sec:limits}

한계 셋이 주장의 범위를 정한다. \ref{sec:pilot}절의 파일럿 증거는 본 설계가 대체한 규칙 변형 아래 셀당 10세션에서 나왔다. 그 횟수는 설계의 동기이지 $\Xstar$의 추정치가 아니다. 거기서 보고한 유일한 유의성 검정(Fisher 정확검정 $p = 0.011$, 10개 중 6개 대 0개)은 시도한 규칙 변형 수로 보정하지 않았다. 눈금은 다섯 단이고, 10라운드 게임에서는 그중 둘이 축 위에서 겹친다. 그래서 50점을 넘는 $\Xstar$는 곡선의 가파르고 드문드문한 구간에서 읽히며, 그 구간의 신뢰구간은 구성상 넓고, 자주 포기하는 모델의 가격을 소수점 한 자리까지 인용하려면 사다리를 먼저 다시 벌려야 한다. 추론 채널 통계는 숙고의 척도로 사고 토큰을 쓰는데, 그것은 길이이지 내용이 아니다. 저장된 추론 텍스트로 내용 분석이 가능하지만 이 글은 하지 않는다~\todo{decide whether a lexicon or judge-based pass on the confidence-call text is in scope}.
